\section{Uppgift 2}\label{sec:uppg02}

\subsection{Instruktioner}
\begin{verbatim}
2. Vilket värde får variablerna x och y efter det att uttrycket (a-g) beräknats
   (varje uttryck ska beräknas enskilt). Lös gärna först uppgiften i huvudet
   (lär man sig mycket på) utan att testa med något program.
   Följande deklaration finns i programmen:

       int x=4;
       int y=2;

       a.  y=x++;
       b.  y=++x;
       c.  y\*=x;
       d.  y+=(x+2);
       e.  y+=--x;
       f.  y+=x++
       g.  x=-y;
\end{verbatim}

\subsection{Lösning}
\subsubsection{Funktion}
% TODO: Eventuell beskrivning av #02 funktionalitet.

\subsubsection{Kommentar}

%                                    värde efter uttryck
% Uttryck        Operation           int x         int y
% =======================================================
% (start)                                4             2
%     (A)        y    =    x++           5             4
%     (B)        y    =  ++x             6             6
%     (C)        y   *=    x             6             36
%     (D)        y   +=   (x + 2)        6             44
%     (E)        y   +=  --x             5             49
%     (F)        y   +=    x++           6             54
%     (G)        x    =   -y;           -54            54
%
%
%
%


\subsubsection{Källkod}
\inputminted[linenos]{java}{src/Lab2Uppg02.java}
\caption{Lab2Uppg02.java}
\label{src:uppg02}
% TODO: Lägg till källkod för #02.


\subsubsection{Skärmdump}
\begin{figure}[htbp]
    \centering
        \includegraphics[width=\linewidth]{img/02.png}
    \caption{Körning av koden till Uppgift \ref{sec:uppg02}}
    \label{fig:uppg02-screenshot}
\end{figure}
% TODO: Lägg till skärmdump av #02.

